\begin{trapbox}
\begin{description}[style=nextline, leftmargin=2.8em, labelsep=0.5em, itemsep=0.35em]
    \item[\textbf{\textcolor{CTrap}{易错点一（方向混淆）}}] 误记不等式方向，认为 $\frac{\ln x + \ln y}{2} \geq \ln\!\left( \frac{x+y}{2} \right)$。
    \item[\textbf{\textcolor{CTrap}{正确理解（凹函数性质）：}}] 由于 $\ln t$ 是凹函数，函数值的平均不超过平均值的函数值，因此不等式方向为 $\leq$。
    \item[\textbf{\textcolor{CTrap}{错因分析（图像误读）：}}] 将凹函数图像误记为“凸向上”，导致方向判断错误。
\end{description}

\medskip
\begin{description}[style=nextline, leftmargin=2.8em, labelsep=0.5em, itemsep=0.35em]
    \item[\textbf{\textcolor{CTrap}{易错点二（忽略定义域）}}] 在应用时忽略 $x, y$ 必须为正实数的条件，错误用于非正数。
    \item[\textbf{\textcolor{CTrap}{正确理解（定义域限制）：}}] 对数函数 $\ln t$ 的定义域是 $(0, +\infty)$，因此 $x, y$ 必须为正。
    \item[\textbf{\textcolor{CTrap}{错因分析（函数未定义）：}}] 忘记对数仅对正数有定义，导致在无效区间使用结论。
\end{description}

\medskip
\begin{description}[style=nextline, leftmargin=2.8em, labelsep=0.5em, itemsep=0.35em]
    \item[\textbf{\textcolor{CTrap}{易错点三（取等条件漏解）}}] 认为等号总是成立，或仅在 $x=y$ 时成立但推广形式中忽略所有变量相等。
    \item[\textbf{\textcolor{CTrap}{正确理解（严格取等）：}}] 等号成立当且仅当所有涉及的变量相等（双变量时 $x=y$，多变量时 $x_1=x_2=\cdots=x_n$）。
    \item[\textbf{\textcolor{CTrap}{错因分析（条件遗漏）：}}] 未注意 Jensen 不等式取等条件要求所有点相同，导致漏解或多解。
\end{description}
\end{trapbox}
